% Original: PDF strana 10, donji slajd; drugi deo.
\slajdid{03-s020b}
\begin{frame}[t,label=03-s020b]{Prenos dozvole signalom EO}
\setlength{\parskip}{0pt}
% element: 03-s020b:e03
Izlaz \textcolor{DEisticanje}{\textbf{EO}} (\emph{enable output})
dozvoljava kodovanje nižem koderu samo ako tekući koder sme da radi,
ali nema aktivan ulaz.
\par\vspace{2mm}
\begin{columns}[c,onlytextwidth]
\begin{column}{.56\textwidth}
% element: 03-s020b:e04
\textcolor{DEisticanje}{\textbf{Invertovani GS nije zamena za EO.}}
\par\vspace{1.5mm}
Za \(EI=0\) uslovljeni GS je nula. Njegova inverzija pogrešno bi
ponovo dozvolila rad nižem koderu.
\par\vspace{1.5mm}
Zabrana mora da se prenosi \textbf{od višeg ka nižem prioritetu},
bez obzira na lokalne zahteve.
\end{column}
\begin{column}{.40\textwidth}\centering
\begin{tikzpicture}[circuit logic US,x=1mm,y=1mm,font=\fontsize{9}{11}\selectfont,
 inv/.style={not gate US,draw,minimum width=8mm,minimum height=7mm},
 gejt/.style={and gate US,draw,logic gate inputs=nn,minimum width=10mm,minimum height=10mm}]
\node[inv] (n) at (11,0) {};
\node[gejt,anchor=input 2] (eo) at (27,0) {};
\draw[DEvod] (n.input)--++(-5,0) node[left] {GS};
\draw[DEvod] (n.output)--(eo.input 2);
\draw[DEvod] (eo.input 1)--++(-5,0)--++(0,11) node[above] {EI};
\draw[DEvod] (eo.output)--++(5,0) node[right] {EO};
\end{tikzpicture}

\par\vspace{1mm}
% element: 03-s020b:d02
\(EO=EI\cdot\overline{GS}\)
\end{column}
\end{columns}
\par\vspace{2mm}
Ni \textbf{neuslovljeni GS} nije dovoljan: to što grupa nema zahtev ne znači
da joj neki viši koder nije već zabranio rad. EO mora čuvati i tu informaciju.
\end{frame}
